\documentclass[10pt,a4paper]{amsart}
\usepackage[margin=19mm]{geometry}
\usepackage{amsmath,amssymb,mathtools}
\usepackage[scr=rsfs,cal=euler]{mathalfa}
\usepackage{xfrac}
\usepackage[unicode,hidelinks,bookmarks=false]{hyperref}
\newcommand{\ie}{i.e.\ }
\newcommand{\cf}{cf.\ }
\newcommand{\resp}{respectively\ }
\newcommand{\Subsection}{\subsection}
\providecommand{\nobreakdash}{\nobreak\hbox{-}}
\setlength{\emergencystretch}{2em}
\multlinegap0pt
\usepackage{amssymb}
\usepackage{braket}
\newtheorem{lemma}{Lemma}
\def\RR{\mathbb{R}}
\title{Irreducibility and regularisation properties of Gaussian quantum Markov semigroups}
\author{Franco Fagnola, Federico Girotti}
\date{Source: arXiv:2609.01547v1 (CC BY 4.0)}
\begin{document}
\maketitle

\noindent\emph{Source and licence.} Irreducibility and regularisation properties of Gaussian quantum Markov semigroups. Version arXiv:2609.01547v1 (CC BY 4.0), distributed by arXiv under the Creative Commons Attribution 4.0 International licence, http://creativecommons.org/licenses/by/4.0/, as recorded in the arXiv licence metadata for this exact version and shown on the abstract page https://arxiv.org/abs/2609.01547v1. Full licence notice: \texttt{licenses/arxiv-2609.01547-CC-BY-4.0.txt}.\par\smallskip

\noindent\textit{Editorial scope.} Counted unit: the article's lemma lem:kerct (source0.tex lines 367--374). The article's complete original proof of it is delivered with it (source0.tex lines 375--390, one \texttt{proof} environment of 16 lines). No mathematics is written, repaired or re-derived: the statement and the proof are the article's own lines, in the article's own notation, and no retained line is substituted or re-flowed.  No further source-local context is needed: every cross-reference of every retained line resolves inside the two delivered spans.  Nothing was omitted from any retained span, so the package carries no omission record.  No retained line contains a citation, so the wrapper carries no bibliography and no bibliography entry.  No retained line contains a figure environment, an image-inclusion command or drawing code, so no image is delivered and the package declares no asset dependency.  Editorial formatting only, in addition: an AMS \texttt{amsart} preamble that restates the article's own declarations and macros used by the retained lines, namely the environments \texttt{lemma} on separate counters, together with the standard packages those lines need.  The retained environments print in this document as Lemma 1 in order of appearance, whereas the article prints the same items with the numbering of its own full document.  The complete native source of the article as distributed in the arXiv e-print for this exact version is bundled unchanged as \texttt{sources/209124/source0.tex}.
\begin{lemma} \label{lem:kerct}
For every $t >0$, the following linear subspaces of $\RR^{2d}$ coincide:
\begin{enumerate}
    \item $\ker(\mathbf{C}_t)$,
    \item the biggest $\mathbf{Z}$-invariant linear subspace in $\ker(\mathbf{C})$,
    \item the biggest $\mathbf{Z}$-invariant linear subspace in $\ker(\mathbf{C}_{Z})$.
\end{enumerate}
\end{lemma}

\begin{proof}
First of all, let us remark that
$$\int_0^t e^{s\mathbf{Z}^*}(\mathbf{Z}^*\mathbf{J} + \mathbf{J}\mathbf{Z})e^{s \mathbf{Z}}ds=e^{t\mathbf{Z}} \mathbf{J} e^{t\mathbf{Z}}-\mathbf{J};$$
therefore for every $\mathbf{z}
\in \RR^{2d}$ one has
$$\langle \mathbf{z}, \int_0^t e^{s\mathbf{Z}^*}(\mathbf{Z}^*\mathbf{J} + \mathbf{J}\mathbf{Z})e^{s \mathbf{Z}}ds \,\mathbf{z} \rangle=\langle e^{t \mathbf{Z}}\mathbf{z}, \mathbf{J} e^{t\mathbf{Z}}\mathbf{z} \rangle -\langle \mathbf{z}, \mathbf{J} \mathbf{z} \rangle=0.$$
Hence, one has that $\ker(\mathbf{C}_t)$ coincide with
$$\ker \left ( \int_0^t e^{s\mathbf{Z}^*}\mathbf{C}_{Z}e^{s \mathbf{Z}}ds\right )\cap \RR^{2d}.$$

In order to conclude it is enough to show that for every positive semidefinite matrix $\mathbf{A}$ and for every matrix $\mathbf{B}$,
$$\ker\left ( \int_0^t e^{s\mathbf{B}^*} \mathbf{A} e^{s \mathbf{B}}ds \right )$$
coincides with the biggest $\mathbf{B}$-invariant subspace in $\ker(\mathbf{A})$; indeed, the statement follows considering $\mathbf{A}=\mathbf{C}$ ($\mathbf{A}=\mathbf{C}_Z$) and $\mathbf{B}=\mathbf{Z}$. This is well known, but we report here the proof for completeness. Let us fix $\mathbf{z} \in \RR^{2d}$, then
\begin{align*}&\langle \mathbf{z},\int_0^t e^{s\mathbf{B}^*} \mathbf{A} e^{s \mathbf{B}}ds \mathbf{z} \rangle=0 \Leftrightarrow  \langle e^{s \mathbf{B}} \mathbf{z},  \mathbf{A} e^{s \mathbf{B}}\mathbf{z} \rangle=0, \, 0\leq s \leq t\\
&\Leftrightarrow e^{s\mathbf{B}} \mathbf{z} \in \ker(\mathbf{A}), \, 0 \leq s \leq t \Leftrightarrow \mathbf{B}^k \mathbf{z} \in \ker(\mathbf{A}), \, k \geq 0.\end{align*}
The second implication is due to the continuity and positivity of $s \mapsto \langle e^{s \mathbf{B}} \mathbf{z},  \mathbf{A} e^{s \mathbf{B}}\mathbf{z} \rangle$, while the last one can be obtained in one sense differentiating and in the other using the Taylor series of the matrix exponential function.
\end{proof}

\end{document}
